\section{PRISM: vectorized probabilistic rounding at variable precision}
\label{sec:prism}

This section describes how \cref{alg:vpsr} is applied to every floating-point
operation of a PyTorch model and what it costs. We implement VPSR in PRISM
(\emph{Probabilistic Rounding with Instruction Set Management}), an open-source
C++17 library that previously applied vectorized SR only at hardware
precision~\citep{gonzalezpepe2026fuzzy}. The Verificarlo compiler
pass~\citep{denis2016verificarlo} replaces each floating-point instruction of
the compiled framework with a call to a PRISM kernel, which computes the exact
result as an EFT pair and rounds it at a virtual precision $t$ taken from a
run-time configuration (\cref{fig:prism-pipeline}). In compiler-instrumented
loops, PRISM SR costs $8.5$--$32.1\times$ native time, $4.4$--$70.5\times$ less
than Verificarlo's MCA random rounding (\cref{sec:prism-bench}).

\begin{figure}[htb]
  \centering
  \begin{tikzpicture}[
    font=\small,
    hbox/.style={draw, rounded corners=2pt, align=center, inner sep=3pt,
        rectangle split, rectangle split parts=2,
        rectangle split part fill={gray!25, gray!6}},
    arr/.style={-{Stealth}, thick}]
    \node[hbox, text width=2.2cm] (op) at (0,0)
    {\textbf{PyTorch}
      \nodepart{two} \code{c = a * b}};
    \node[hbox, text width=2.9cm] (entry) at (4.16,0)
    {\textbf{PRISM call}
      \nodepart{two} \code{mulf32x16(a,b)}};
    \node[hbox, text width=2.8cm] (kernel) at (8.43,0)
    {\textbf{VPSR kernel}
      \nodepart{two} EFT $(\sigma,\tau)$,\\ then \cref{alg:vpsr}};
    \node[hbox, text width=2.2cm] (out) at (11.95,0)
    {\textbf{Result}
      \nodepart{two} $c$ rounded at\\ precision $t$};
    \node[hbox] (cfg) at (8.43,1.7)
    {\textbf{Configuration}
      \nodepart{two} virtual precision $t$};
    % \node[hbox, text width=2.9cm] (hook) at (4.16,2.3)
    % {\textbf{Module hooks}
    %   \nodepart{two} \cref{sec:exp-scoping}};
    \draw[arr] (op) -- node[above, font=\scriptsize, align=center]
    {Verificarlo\\ pass} (entry);
    \draw[arr] (entry) -- node[below, font=\scriptsize] {16 lanes} (kernel);
    \draw[arr] (kernel) -- (out);
    % \draw[arr] (hook) -- node[above, font=\scriptsize] {run time} (cfg);
    \draw[arr] (cfg) -- (kernel);
  \end{tikzpicture}
  \caption{Execution of one instrumented operation. At compile time, the
    Verificarlo pass replaces a 512-bit \code{binary32} multiplication with a
    call to the PRISM entry point for $16$ lanes. At run time, the kernel
    rounds all lanes with \cref{alg:vpsr} at the precision set.}
  \label{fig:prism-pipeline}
\end{figure}

\subsection{Instrumentation and kernels}
\label{sec:prism-kernels}

The Verificarlo pass maps each scalar or vector floating-point instruction to a
PRISM entry point for the same operation, format and lane count, so PyTorch's
source is unchanged. Each call rounds all its lanes in one kernel invocation.
We use PRISM's static build, which fixes the instruction set at compile time.

Each kernel computes the EFT pair $(\sigma,\tau)$ of \cref{sec:background-eft}
and then applies \cref{alg:vpsr}. SR and RN share the kernel: SR draws $z$ from
$r$ random bits, and RN fixes $z=1/2$. Because the final comparison of
\cref{alg:vpsr} is non-strict, an exact midpoint rounds away from
$\tilde\sigma$, and hence away from zero; giving \emph{untied} RN.
FMA is rounded as one operation, using the \code{ErrFmaNearest}
transform~\citep{boldo2011exact} to represent $ab+c$ exactly.

The kernels are branch-free. Every data-dependent condition, in
\cref{alg:vpsr} and in the EFTs, is computed as a lane mask and resolved by
selection, so all lanes follow one instruction stream whatever their values,
including infinities and NaNs. The kernels are written once against Google
Highway~\citep{highway}, a C++ SIMD abstraction, and compile to its x86, Arm,
RISC-V and WebAssembly targets; where no hardware FMA exists, \code{TwoProdFMA}
uses the emulated FMA of \citet{graillat2023emulation}.

\subsection{Run-time configuration}
\label{sec:prism-config}

The virtual precisions for \code{binary32} and \code{binary64}, the rounding
mode and the generator state are run-time settings, so changing them requires
no recompilation. In our experiments, PyTorch module hooks change them before
and after each target module (\cref{sec:exp-scoping}). Kernels read the
settings from a per-thread cache. A process-wide setter publishes new settings
and advances a configuration epoch, and each worker thread reloads its cache
when it observes a new epoch on its next operation, so a change made from the
Python thread reaches all of PyTorch's worker threads.

\subsection{Performance}
\label{sec:prism-bench}

\paragraph{Cost model.}
\label{sec:prism-cost}
Emulation replaces one instruction with an EFT (two to six operations), about
twenty-five vector operations for the rounding decision, and part of a random
draw. The control flow is fixed, so this cost does not depend on the operands.
The slowdown relative to native execution is therefore smaller when native
execution is memory-bound and larger when a call covers few lanes.

\paragraph{Benchmark.}
We time elementwise addition, multiplication, division and FMA over arrays of
$65{,}536$ values in \code{binary32} and \code{binary64}. Each loop is compiled
with \code{-O3} for 512-bit AVX-512 vectors in three builds: without
instrumentation (the native baseline), with the PRISM SR pass, and with
Verificarlo's MCA QUAD backend in random-rounding (RR) mode
(\cref{sec:background-eft}). We repeat the measurement $30$ times on Intel Xeon
Gold 6130 @ 2.1~GHz CPU (cache sizes: 32~KB L1, 1~MB L2, 22~MB L3 shared).

\begin{table}[t]
  \centering
  \caption{Slowdown of PRISM SR versus MCA RR relative to native execution
  ($N=65{,}536$, 30 runs; $t=24$ for \code{binary32}, $t=53$ for
  \code{binary64}). PRISM and MCA entries are the median ratios to
  native time for the same operation and data type (lower is better); MCA/PRISM
  is the median per-run ratio of MCA to PRISM time.}
  \label{tab:prism-vector-perf}
  \begin{tabular}{lcccccc}
    \toprule
                                                   & \multicolumn{3}{c}{\code{binary32}} &
    \multicolumn{3}{c}{\code{binary64}}                                                                                                                    \\
    \cmidrule(lr){2-4}\cmidrule(lr){5-7} Operation & PRISM                               & MCA
                                                   & MCA/PRISM                           & PRISM             & MCA   & MCA/PRISM                           \\
    \midrule
    add                                            & $31.1$                              & $142.4$           & $4.7$ & $11.6$           & $115.6$ & $10.0$ \\
    \rowcolor{black!6} mul                         & $32.1$                              & $143.7$           & $4.5$ & $11.3$           & $163.1$ & $14.6$ \\
    div                                            & $25.6$                              & $112.7$           & $4.4$ & $12.0$           & $229.9$ & $19.2$ \\
    \rowcolor{black!6} fma                         & $19.1$                              & $\phantom{0}88.3$ & $4.7$ & $\phantom{0}8.5$ & $597.4$ & $70.5$ \\
    \bottomrule
  \end{tabular}
\end{table}

\paragraph{Results.}
PRISM SR incurs an overhead of $\times8.5$--$32.1$, whereas MCA RR incurs
$\times88$--$597$ (\cref{tab:prism-vector-perf}). Compared to MCA RR, PRISM SR
is up to $\times4.7$ faster than for \code{binary32} operations and up to
$\times70.5$ faster for \code{binary64}. The gap comes from how each backend
obtains the exact result it rounds. MCA RR computes the result in a wider format
and adds its random perturbation there: in \code{binary64} for \code{binary32}
operations, which the hardware supports, and in \code{binary128} for
\code{binary64} operations, which most CPUs emulate in software. PRISM recovers
the exact result with EFTs in the working format, so both formats use hardware
instructions only. Performance is nearly independent of the virtual precision:
across precisions, PRISM times vary by less than $1\%$ and MCA times by at most
$5\%$.
